\summarychapter{18}{The Search for Unity}{Charlemagne and Christendom}

\begin{inbrief}
The Middle Ages knew neither the secular sovereign state nor the church as a voluntary society.
Drawing on Augustine's \emph{City of God}, Charlemagne grafted the universal church onto the Roman
idea of empire and produced Christendom, one society blending eternal and temporal concerns. The
chapter traces the rise of the Franks, the alliance with the papacy, the decline into feudalism and
the reform that set pope against emperor.
\question{How could the kingdom that Jesus said was ``not of this world'' become so much a part of
worldly power?}
\end{inbrief}

On April 25, 799, armed men attacked Pope Leo III as he led a procession through Rome, pulled him from
his horse and carried him off. Supporters of the previous pope, Adrian I, accused him of perjury and
adultery. Rescued by night, Leo appealed to the papacy's traditional protector, the king of the
Franks. Charles crossed the Alps and presided over a great assembly. On December 23, 800, the pope
cleared himself by oath, and on Christmas Day, as Charles knelt at St.~Peter's, Leo set a golden crown
on his head while the congregation hailed him as Augustus, ``crowned by God.'' The Christian Roman
Empire was restored. The answer to how it happened, Shelley says, lies in the persistence of an idea
and the rise of a powerful kingdom.

\ssection{Ideas Die Slowly}
Long after its fall the Roman Empire ruled the imagination. Amid many warring kingdoms people longed
for the old unity and expected a new Roman Empire. The Franks seemed destined to restore it. Clovis,
with the support of the Catholic Church, had made them dominant, but his Merovingian dynasty decayed
through the Germanic custom of dividing the kingdom among the king's sons, with civil war, murder and
intrigue as the result. Power passed to the landed aristocrats and above all to the ``mayor of the
palace.'' In 714 Charles Martel, Charlemagne's grandfather, took that office and left the kings as
figureheads. Muslims from North Africa had invaded Spain in 711 and destroyed the Visigothic kingdom
by 718. In 732 Charles met their raiders near Tours and drove them back, which earned him the name
Martel, ``the Hammer''; they never again seriously threatened central Europe.

His son Pepin the Short (741--768) wanted the title that matched his power and asked the pope to rule
that whoever held actual power should be the legal ruler. With papal blessing Boniface crowned him
king in 751, the last Merovingian was put away in a monastery, and three years later the pope crossed
the Alps to anoint Pepin in the Old Testament manner. The pope needed a protector. The Lombards had
taken Ravenna, seat of Byzantine government in Italy, and threatened Rome. In 756 a Frankish army
forced the Lombards to give up their conquests, and Pepin conferred the territory on the pope. This
``Donation of Pepin'' made the pope temporal ruler of the Papal States, a strip across Italy from
coast to coast. The alliance had lasting effects. It hastened the separation of the Latin from the
Greek church by giving the papacy a western ally in place of the Byzantines; it created the Papal
States, a factor in Italian politics until the late nineteenth century; and the anointing gave western
kingship a religious sanction that fed the later rivalry of pope and emperor. When Charles succeeded
in 768 he sought military power to crush his enemies, religious power to direct souls and
intellectual power to instruct minds. His success made Europe nominally Christian for a thousand
years.

\ssection{Architect of an Empire}
His biographer Einhard pictures a tall, strong horseman, a warrior-king who campaigned every year and
also wanted effective government. He added four regions by force. In Spain, after a first crossing of
the Pyrenees in 778, he pushed the Muslims back to the Ebro and set up the Spanish Mark around
Barcelona. He conquered the Bavarians and, in thirty-two campaigns, the pagan Saxons between Rhine
and Elbe, dividing Saxony into bishoprics, building monasteries and punishing with death such
offenses as eating meat in Lent, cremating the dead or pretending to be baptized; his brutality was
extreme by any measure. In six campaigns he crushed the Avars and founded the East Mark on the Danube,
later Austria. And in Italy, at the pope's request, he defeated the Lombards in 774, made himself
their king and confirmed the Donation of Pepin.

That was the prelude to 800. The pope needed protection and Charlemagne divine sanction. He claimed
not to have known of the plan, but he lived up to the title, styled himself Roman Emperor and
Augustus, had all officers swear to him as Caesar, and in 812 won the recognition of Constantinople.
The coronation showed how alive the memory of Rome was; it also began a long struggle. In theory
church and state were two aspects of one Christendom, one securing spiritual blessings, the other
justice and welfare. In fact pope and emperor were rivals. The papal party argued that time depends
on eternity and so the emperor on the pope; the imperial party, that God ordained the Christian state
too, and that an emperor may even correct a negligent pope. In his lifetime Charlemagne left no doubt
who ruled. His \emph{missi dominici}, pairs of envoys, usually a bishop and a lay noble, inspected
local government everywhere. He fostered a revival of learning, decreeing in 789 that every monastery
keep a school and founding a palace school at Aix-la-Chapelle under the Anglo-Saxon Alcuin, who wrote
textbooks on grammar, spelling, rhetoric and logic. From a new center in the north he extended
Christian civilization, restored a measure of order after three centuries and revived an imperial
ideal that lasted until Napoleon ended the Holy Roman Empire in 1806.

\ssection{The Decline into Feudalism}
The empire was only a breathing space. Under weak successors it broke up in civil wars and new
invasions, and as the Vikings came, people gave their lands and themselves to local lords for
protection. Feudalism was government exercised locally by private persons. It fused three elements:
the personal bond of lord and vassal, the fief (usually land) given to the vassal, and the private
exercise of government over vassals and fiefs. In the act of homage the vassal knelt and became the
lord's ``man''; in the oath of fealty he swore loyalty; in investiture he received a symbol of his
jurisdiction. The contract was sacred. The lord owed protection and justice, the vassal chiefly
military service, forty days a year.

\begin{sidebox}{Praying Hands}
By the end of the eleventh century the vassal knelt, joined his hands and held them out; the lord
clasped them, raised him and kissed him. Hands joined and extended have become the posture of prayer.
\end{sidebox}

The church was drawn in. Threatened by Vikings and Magyars, bishops and abbots became vassals of the
barons in France and the kings in Germany, holding fiefs and owing service, which strained their
loyalty to the pope. In the tenth and early eleventh centuries the papacy itself was a prize of the
Roman nobles. On the positive side the church tried to add Christian virtues to chivalry and to limit
warfare through the Peace of God, which barred from the sacraments those who plundered holy places or
harmed noncombatants, and the Truce of God, which closed certain days and seasons to fighting; both
were largely ineffective. Some unity returned when the German king Otto the Great revived the empire
in 962, but with it came the old rivalry, for the German kings interfered even in papal elections.
The quarrel centered on lay investiture. In theory a bishop received spiritual authority from the
church and civil authority from the king; in practice kings and lords controlled both, above all in
Germany, where the church was in essence a state church.

\ssection{Setting in Order the House of God}
The church first had to reform itself. The impulse came from the Benedictine house of Cluny in
Burgundy, founded in 910. Its program began with the monasteries and went on to demand clerical
celibacy and the end of simony, the buying and selling of church office (named after Simon the
magician in Acts 8). Its final aim was to free the whole church from secular control and place it
under the pope. Some three hundred Cluniac houses were freed from lay control, and in 1059 the College
of Cardinals was created to elect the popes. The man behind this was the archdeacon Hildebrand. As
Gregory VII (1073--1085) he aimed at a Christian commonwealth under papal control, with the spiritual
power supreme over the temporal. In 1075 he forbade lay investiture on pain of excommunication,
virtually declaring war on Europe's rulers.

The clash came with the emperor Henry IV, who had installed his own archbishop of Milan. Summoned to
Rome, Henry had a synod of German bishops declare Gregory a usurper in 1076. Gregory excommunicated
and deposed him and released his subjects from their oaths. Faced with a revolt of his nobles, Henry
came in January 1077 to Canossa and stood three days barefoot in the snow as a penitent until the pope
absolved him. The humiliation did not end the quarrel, and contemporaries made little of it, but the
papacy had gained. The Concordat of Worms settled investiture in 1122: the church elects, in the
presence of the emperor or his representative. Later popes added little to Gregory's theory that
Christian society is organized under the pope, who is kept from error by Peter's presence in his
successors. Many of these claims now seem intolerable, yet Shelley holds that they stood for two
principles no Christian can contest: the spiritual has primacy over the secular in human loyalties,
and humanity finds true unity only in Christ. In the Middle Ages Europe became conscious of itself as
a unity, and the church reached a power over people's lives, for the most part used helpfully, that
it has never had since.

\begin{keyfacts}
\keypeople{Charlemagne; Pope Leo III; Charles Martel; Pepin the Short; Boniface; Einhard; Alcuin; Otto
the Great; Hildebrand (Gregory VII); Henry IV.}
\keydates{%
  \kd{732}{Charles Martel defeats the Muslims near Tours}
  \kd{751}{Pepin the Short crowned king of the Franks}
  \kd{756}{Donation of Pepin: origin of the Papal States}
  \kd{768}{Charlemagne becomes king of the Franks}
  \kd{800}{Charlemagne crowned emperor by Leo III on Christmas Day}
  \kd{910}{Monastery of Cluny founded}
  \kd{962}{Otto the Great revives the empire in the West}
  \kd{1059}{College of Cardinals created to elect the pope}
  \kd{1073--1085}{Pontificate of Gregory VII}
  \kd{1077}{Henry IV does penance at Canossa}
  \kd{1122}{Concordat of Worms settles the investiture dispute}}
\keyterms{Christendom; mayor of the palace; Donation of Pepin; Papal States; \emph{missi dominici};
feudalism; vassal, fief, homage, investiture; Peace and Truce of God; Cluny; simony; lay investiture;
College of Cardinals.}
\end{keyfacts}
